\chapter{SYSTEM ARCHITECTURE}

\section{OVERVIEW}

Stockify employs a modern three-tier architecture designed for scalability, security, and maintainability. The system follows clear separation of concerns with distinct presentation, application, and data layers.

\section{ARCHITECTURAL DESIGN}

\subsection{Three-Tier Architecture}
\begin{itemize}
    \item \textbf{Presentation Layer:} React-based frontend with responsive UI
    \item \textbf{Application Layer:} Node.js/Express.js REST API server  
    \item \textbf{Data Layer:} MongoDB Atlas database with document storage
\end{itemize}

\subsection{Design Principles}
\begin{itemize}
    \item \textbf{Separation of Concerns:} Clear boundaries between layers
    \item \textbf{Multi-Tenancy:} Secure data isolation for different users
    \item \textbf{RESTful Design:} Stateless, resource-oriented API
    \item \textbf{Scalability:} Horizontal and vertical scaling capabilities
\end{itemize}

\section{FRONTEND ARCHITECTURE}

\subsection{Component Hierarchy}
The React frontend follows a hierarchical structure with Pages (route components), Layouts (shared navigation), Components (reusable UI elements), and Hooks (custom state management).

\subsection{State Management}
State management utilizes React Context for global authentication state, local component state for UI interactions, custom hooks for data fetching, and URL-based state for navigation.

\subsection{Data Flow}
User interactions trigger component events, which call API services, make HTTP requests to the backend, update component state, and re-render the UI with new data.

\section{BACKEND ARCHITECTURE}

\subsection{MVC Pattern}
The backend implements Model-View-Controller architecture with Mongoose models for data structure, controllers for business logic, routes for API endpoints, and services for shared functionality.

\subsection{Middleware Stack}
Express.js middleware handles authentication (JWT verification), validation (Joi schemas), error handling, logging, and security (CORS, rate limiting).

\subsection{API Design}
RESTful API follows resource-based URLs (/api/products, /api/sales), proper HTTP methods and status codes, consistent JSON format, and built-in pagination for large datasets.

\section{DATABASE ARCHITECTURE}

\subsection{Document-Oriented Design}
MongoDB provides flexible schema design, embedded documents for performance, rich query capabilities with aggregation pipelines, and horizontal scaling with built-in sharding.

\subsection{Multi-Tenant Isolation}
Data isolation uses user scoping with createdBy fields, automatic filtering middleware, complete privacy between users, and efficient indexing on tenant identifiers.

\subsection{Key Collections}
Main collections include Users (profiles and authentication), Products (catalog with pricing), Sales (transactions with line items), Categories (organization), Customers (information), and Suppliers (vendor management).

\section{SECURITY ARCHITECTURE}

\subsection{Authentication}
Multi-provider authentication supports JWT tokens for stateless authentication, Google OAuth integration, bcrypt password hashing, and automatic token refresh.

\subsection{Authorization}
Access control implements role-based permissions, resource ownership validation, API endpoint protection, and comprehensive input validation.

\section{PERFORMANCE OPTIMIZATION}

\subsection{Frontend Performance}
Optimization includes code splitting and lazy loading, bundle optimization with tree shaking, browser caching for static assets, and debounced search functionality.

\subsection{Backend Performance}
Server optimization uses strategic database indexing, connection pooling, response caching for frequent queries, and efficient pagination for large datasets.

\section{DEPLOYMENT ARCHITECTURE}

\subsection{Development Environment}
Local development uses Vite development server (port 5173), Node.js with nodemon (port 5000), MongoDB Atlas connection, and hot reload for efficiency.

\subsection{Production Environment}
Production deployment includes static build served by web server, Node.js production server with PM2, MongoDB Atlas with automated backups, and HTTPS encryption.

\section{MONITORING AND INTEGRATION}

\subsection{System Monitoring}
Monitoring strategy includes centralized error logging, performance metrics tracking, user analytics, and automated health checks.

\subsection{External Integrations}
Third-party services include Google OAuth authentication, email notification systems, file storage solutions, and analytics tools.

\section{SCALABILITY CONSIDERATIONS}

The architecture supports future growth through stateless API design for server replication, MongoDB Atlas automatic scaling, CDN integration for static assets, and load balancing capabilities.

\section{CONCLUSION}

The Stockify architecture provides a solid foundation for scalable, secure, and maintainable inventory management. The modular design enables future enhancements while maintaining system stability and performance.